\unitopen{39}
\begin{candos}
\item zreferować wypowiedź innej osoby w mowie zależnej z \de{Konjunktiv I};
\item znać formy \de{Konjunktiv I} (\de{sei, habe, werde, komme}) i kiedy zastępuje je \de{Konjunktiv II};
\item odróżnić referowanie neutralne od dystansującego i użyć \de{laut}, \de{angeblich}, \de{offenbar};
\item streścić artykuł lub spotkanie w trzeciej osobie.
\end{candos}

\section{Słownictwo}
\vocabtable{39}

\section{Po co \de{Konjunktiv I}}
Gdy cytujesz, co ktoś powiedział, używasz \emph{mowy zależnej}. W mowie potocznej wystarcza Indikativ (\de{Er sagt, er hat keine Zeit}). W piśmie
i w mediach używa się \emph{Konjunktiv I}, żeby zaznaczyć, że \emph{nie bierzesz odpowiedzialności} za treść:
\de{Der Sprecher sagte, das Projekt \textbf{sei} fertig.}

\begin{center}\small
\begin{tabular}{@{}lll@{}}
\toprule
\textsf{\textbf{Bezokolicznik}} & \textsf{\textbf{er/sie/es (Konj. I)}} & \textsf{\textbf{ich, wir, sie (mnoga)}}\\
\midrule
\de{sein} & \de{sei} & \de{ich sei}, \de{wir seien}\\
\de{haben} & \de{habe} & \de{ich habe} (= Indik.) $\rightarrow$ \de{hätte}\\
\de{werden} & \de{werde} & \de{wir werden} (= Indik.) $\rightarrow$ \de{würden}\\
\de{kommen} & \de{komme} & \de{wir kommen} (= Indik.) $\rightarrow$ \de{kämen}\\
\de{können} & \de{könne} & \de{sie können} (= Indik.) $\rightarrow$ \de{könnten}\\
\bottomrule
\end{tabular}
\end{center}
Reguła: gdy forma \de{Konjunktiv I} jest taka sama jak Indikativ, zastępujesz ją \de{Konjunktiv II} (jednostka~29), żeby było widać, że to mowa zależna.
Przeszłość: \de{habe}/\de{sei} + Partizip II: \de{Er sagte, er \textbf{habe} das Projekt \textbf{abgeschlossen}.}

\begin{zasada}[Trzy poziomy dystansu]
(1) \de{Er sagt, dass er krank ist.} (neutralnie, mowa potoczna) (2) \de{Er sagt, er sei krank.} (formalnie, dystans) (3) \de{Er behauptet,
er sei krank.} lub \de{Er ist angeblich krank.} (wątpliwość). Wybór sygnalizuje stosunek do informacji.
\end{zasada}

\begin{falle}
W polskim mowa zależna zachowuje czas wypowiedzi i nie ma osobnego trybu; po niemiecku Konjunktiv I jest typowy dla prasy i
nie powinien dominować w mowie. Błąd: użycie Konjunktiv I w rozmowie z kolegą (brzmi sztucznie). Drugi: stosowanie formy identycznej
z Indikativ bez zamiany na Konjunktiv II (\de{Sie sagten, sie haben keine Zeit} w piśmie wymaga \de{hätten}).
\end{falle}

\begin{szwajcaria}
Szwajcarskie media (\de{SRF}, \de{NZZ}) używają Konjunktiv I w relacjach: \de{Die Behörden teilten mit, es gebe keine Verletzten.}
Do mowy codziennej wystarczy Ci Indikativ i \de{dass}-zdanie. Ćwicz rozumienie, bo Konjunktiv I jest wskazówką, że autor się dystansuje.
\end{szwajcaria}

\begin{praca}
Protokół: \de{Die Kundin erklärte, sie sei mit der Lösung zufrieden und habe keine weiteren Fragen.} \de{Der Sprecher des Anbieters
sagte, der Dienst werde am Freitag wieder verfügbar sein.} Zwróć uwagę na dystans: autor nie potwierdza, że to prawda.
\end{praca}

\section{Dialogi}
\textbf{Dialog 1 --- relacja z rozmowy}
\begin{dialog}
\turn{A}{Was hat der Kunde gesagt?}{Co powiedział klient?}
\turn{B}{Er sagte, er sei mit dem Angebot zufrieden, aber er habe noch Fragen zum Preis.}{Powiedział, że jest zadowolony z oferty, ale ma jeszcze pytania o cenę.}
\turn{A}{Hat er etwas zum Termin gesagt?}{Czy coś powiedział o terminie?}
\turn{B}{Er meinte, er könne am Dienstag kommen.}{Stwierdził, że może przyjść we wtorek.}
\end{dialog}

\section{Ćwiczenia}
\begin{cwiczenia}
\exitem{Konjunktiv I: \de{sein} (er).}{\de{sei}}
\exitem{Konjunktiv I: \de{haben} (er).}{\de{habe}}
\exitem{Konjunktiv I: \de{werden} (er).}{\de{werde}}
\exitem{Konjunktiv I: \de{können} (er).}{\de{könne}}
\exitem{Mowa zależna: \de{Er sagt: „Ich bin krank.”}}{\de{Er sagt, er sei krank.}}
\exitem{Mowa zależna: \de{Sie sagte: „Wir haben keine Zeit.”}}{\de{Sie sagte, sie hätten keine Zeit.} (\de{haben} = Indik., więc Konj. II)}
\exitem{Mowa zależna, przeszłość: \de{Er sagte: „Ich habe das Projekt abgeschlossen.”}}{\de{Er sagte, er habe das Projekt abgeschlossen.}}
\exitem{Użyj \de{laut} i \de{Bericht}: „Według raportu koszty rosną.”}{\de{Laut dem Bericht steigen die Kosten.}}
\exitem{Wyraź wątpliwość: \de{Er ist angeblich krank.} $\rightarrow$ po polsku.}{Rzekomo jest chory.}
\exitem{Przetłumacz: „Rzecznik powiedział, że usługa będzie dostępna w piątek.”}{\de{Der Sprecher sagte, der Dienst werde am Freitag verfügbar sein.}}
\end{cwiczenia}

\weekend{Znajdź krótką wiadomość SRF lub DW i napisz jej streszczenie w mowie zależnej (osiem zdań). Potem przekształć je
na język potoczny z \de{dass}. Nagraj oba warianty.}
